%!TEX root = l2kurz.tex
% bilder.tex - Teil der LaTeX2e-Kurzbeschreibung
% Siehe https://github.com/texdoc/l2kurz

\section{Setzen von Bildern} \label{graphics}

Lädt man im Vorspann des Dokuments das Paket \texttt{graphicx}, dann kann man
Bilder, die mit einem anderen Programm erzeugt wurden,  mit dem Befehl
\lstinline|\includegraphics| in das Dokument einfügen. Der Dateiname ist als
Argument des Befehls anzugeben. Welche Dateiformate verarbeitet werden können,
hängt vom jeweiligen \TeX-System ab. Das häufig benutzte PDF\TeX{} kann
Grafiken in den Formaten PDF, PNG JPEG und JBIG2 verarbeiten.

% http://openclipart.org/detail/104863/ocean-theme-papercut-by-last-dino
% Public Domain (CC0)
% http://creativecommons.org/publicdomain/zero/1.0/

\begin{LTXexample}
Hier \includegraphics
       [width=3cm]{ozean}
ist ein Bild.
\end{LTXexample}

\noindent Wird das Paket \texttt{graphicx} mit der Option \texttt{[draft]} geladen,
dann erscheint anstelle des Bildes nur ein Rahmen entsprechend
der tatsächlichen Bildgröße mit dem Namen des Grafikfiles, 
was die Bearbeitung beschleunigt und für Probeausdrucke nützlich ist.

Weitere Informationen zum Einbinden von Bildern finden Sie in der
Online"=Dokumentation \cite{grfguide}, im \textit{Graphics Companion}
\cite{grfcomp} und in K.~Reckdahls empfehlenswertem  Tutorium \cite{epslatex}.



